\scanpage{5}

\deleted{Utilisant le corollaire 2 \deleted{lorsque $S=\mathrm{Spec}(k\unsure{[\epsilon]|[t]|(\varepsilon)})$} pour $S'=\mathrm{Spec}(\underline D(\underline{\unsure{I|J}}))$ où $\underline I$ est un faisceau coh. sur $S$, $\underline D(\underline I)=\underline{\mathcal O}_S+\underline I$ ($\underline I$ de carré nul) et tous les $S'$-automorphismes de $C'=C\times_S S'$ qui induisent l'identité sur $S$, on trouve}

\deleted{\unsure{Prop.}} \deleted{\unsure{Partons}} \uline{Corollaire 3} Soient $C,C'$ deux courbes sur $S$ satisfaisant les hyp. du th. Alors
\[
\underline{\underline{\mathrm{Is}}}_S(C,C')\longrightarrow\underline{\underline{\mathrm{Is}}}_S(J^0_{C/S},J^0_{C'/S})\quad\text{est une}
\]
immersion fermée.

Soit maintenant $A$ un \uline{schéma abélien} \uline{polarisé} sur $S$ \inserted{de dim relative $g$} (polarisé au sens strict : on se donne une section du schéma de \unsure{Néron-Sévéri|Néron-Séveri|Nér.-Sév.}). Pour \unsure{tt} $\deleted{S}T'$ sur $S$, on désigne par $F(T)$ l'ensemble des classes, à isom près, de courbes $C$ sur $\deleted{\unsure{S'|S_T}}T'$, satisfaisant aux conditions générales de ce \unsure{N°|Th.}, et \uline{munies} d'un isom de $J^0_{C/T}$ sur $A$ \deleted{\unsure{tel de façon que}} \inserted{compatible avec} les polarisations.

% archive sheet number: 4
